\documentclass[11pt,a4paper]{article}
\usepackage[utf8]{inputenc}
\usepackage[T1]{fontenc}
\usepackage[french]{babel}
\usepackage[a4paper,top=1.75cm,bottom=1.9cm,left=1.85cm,right=1.85cm]{geometry}
\usepackage{amsmath,amssymb}
\usepackage{array}
\usepackage{booktabs}
\usepackage{enumitem}
\usepackage{eurosym}
\usepackage{fancyhdr}
\usepackage{hyperref}
\usepackage{lmodern}
\usepackage{inconsolata}
\usepackage{listings}
\usepackage{microtype}
\usepackage{tabularx}
\usepackage[most]{tcolorbox}
\usepackage{titlesec}
\usepackage{xcolor}

\renewcommand{\familydefault}{\sfdefault}

\definecolor{ink}{HTML}{202124}
\definecolor{muted}{HTML}{5F6368}
\definecolor{line}{HTML}{DADCE0}
\definecolor{paper}{HTML}{F8F9FA}
\definecolor{panel}{HTML}{FFFFFF}
\definecolor{accent}{HTML}{8A1538}
\definecolor{accentsoft}{HTML}{F8EEF2}
\definecolor{green}{HTML}{0B6B57}
\definecolor{greensoft}{HTML}{E6F4EF}
\definecolor{amber}{HTML}{9A5B00}
\definecolor{ambersoft}{HTML}{FFF7E6}
\definecolor{blue}{HTML}{245B8E}
\definecolor{bluesoft}{HTML}{EAF2FA}
\definecolor{rose}{HTML}{A84D77}
\definecolor{rosesoft}{HTML}{FAF0F5}
\definecolor{codeback}{HTML}{F8F9FA}
\definecolor{codemark}{HTML}{DADCE0}

\hypersetup{
  colorlinks=true,
  linkcolor=accent,
  urlcolor=green,
  pdfborder={0 0 0}
}

\setlength{\parindent}{0pt}
\setlength{\parskip}{5pt}
\setlist[itemize]{leftmargin=1.35em,itemsep=2pt,topsep=2pt}
\setlist[enumerate]{leftmargin=1.55em,itemsep=3pt,topsep=3pt}
\setlist[enumerate,1]{label=\textbf{\alph*.}}

\pagestyle{fancy}
\fancyhf{}
\fancyhead[L]{\sffamily\footnotesize Python pour économistes}
\fancyhead[R]{\sffamily\footnotesize Université de Lille -- L2 Économie}
\fancyfoot[C]{\sffamily\footnotesize\color{muted}\thepage}
\renewcommand{\headrulewidth}{0.2pt}
\renewcommand{\headrule}{\hbox to\headwidth{\color{line}\leaders\hrule height \headrulewidth\hfill}}
\setlength{\headheight}{22pt}

\titleformat{\section}
  {\sffamily\Large\bfseries\color{ink}}
  {}{0pt}{}
  [\vspace{-0.25em}{\color{accent}\titlerule[0.7pt]}]
\titleformat{\subsection}
  {\sffamily\large\bfseries\color{ink}}
  {}{0pt}{}

\lstdefinelanguage{terminal}{
  morekeywords={ls,cd,pwd,mkdir,python,python3,code,dir,type,echo},
  sensitive=true
}

\lstset{
  basicstyle=\ttfamily\footnotesize,
  backgroundcolor=\color{codeback},
  frame=single,
  framesep=5pt,
  framerule=0.25pt,
  rulecolor=\color{codemark},
  language=Python,
  keywordstyle=\color{accent}\bfseries,
  commentstyle=\color{muted},
  stringstyle=\color{green},
  numberstyle=\scriptsize\color{muted},
  numbers=left,
  numbersep=8pt,
  showstringspaces=false,
  breaklines=true,
  aboveskip=0.7em,
  belowskip=0.7em,
  xleftmargin=1.1em,
  framexleftmargin=1em,
  columns=fullflexible,
  keepspaces=true,
  upquote=true,
  literate=
   {à}{{\`a}}1 {â}{{\^a}}1 {ä}{{\"a}}1
   {ç}{{\c c}}1
   {é}{{\'e}}1 {è}{{\`e}}1 {ê}{{\^e}}1 {ë}{{\"e}}1
   {î}{{\^\i}}1 {ï}{{\"\i}}1
   {ô}{{\^o}}1 {ö}{{\"o}}1
   {ù}{{\`u}}1 {û}{{\^u}}1 {ü}{{\"u}}1
   {À}{{\`A}}1 {Â}{{\^A}}1 {Ä}{{\"A}}1
   {Ç}{{\c C}}1
   {É}{{\'E}}1 {È}{{\`E}}1 {Ê}{{\^E}}1 {Ë}{{\"E}}1
   {Î}{{\^I}}1 {Ï}{{\"I}}1
   {Ô}{{\^O}}1 {Ö}{{\"O}}1
   {Ù}{{\`U}}1 {Û}{{\^U}}1 {Ü}{{\"U}}1
   {œ}{{\oe}}1 {Œ}{{\OE}}1
   {–}{{--}}1 {—}{{---}}1
}

\newcommand{\code}[1]{\texttt{#1}}
\newcommand{\pp}{\,points de pourcentage}
\newcommand{\taglabel}[1]{\scriptsize\sffamily\bfseries #1}

\newcounter{exercise}
\newcounter{extension}
\newcounter{logic}

\newcommand{\workbooktitle}[4]{%
  \setcounter{exercise}{0}%
  \setcounter{extension}{0}%
  \setcounter{logic}{0}%
  \fancyhead[L]{\sffamily\footnotesize #1 -- #2}%
  \begingroup\raggedright
  {\sffamily\Large\bfseries\color{ink}#1 -- #2\par}
  \vspace{0.25em}
  {\sffamily\small\color{muted}Python pour économistes -- Université de Lille, L2 Économie\par}
  \vspace{0.85em}
  {\sffamily\bfseries\color{accent}Question fil rouge.} #3\par
  \vspace{0.35em}
  {\sffamily\bfseries\color{green}Livrables.} #4\par
  \vspace{0.9em}
  {\color{line}\hrule height 0.45pt}
  \vspace{0.85em}
  \endgroup
}

\newenvironment{objectivebox}[1]{%
  \subsection*{#1}
  \vspace{-0.55em}
}{\vspace{0.15em}}

\newenvironment{routebox}[1]{%
  \subsection*{#1}
  \vspace{-0.45em}
}{\vspace{0.15em}}

\newenvironment{deliverablebox}[1]{%
  \subsection*{#1}
  \vspace{-0.45em}
}{\vspace{0.15em}}

\newenvironment{methodbox}[1]{%
  \subsection*{#1}
  \vspace{-0.45em}
}{\vspace{0.15em}}

\newenvironment{pathwaybox}[1]{%
  \par\addvspace{0.8em}
  {\sffamily\bfseries\color{blue}#1\par}
  \vspace{-0.25em}
  \begingroup\leftskip=0.85em\relax
  \noindent\color{ink}
}{\par\endgroup\addvspace{0.25em}}

\newenvironment{pseudocodeexercise}[1]{%
  \refstepcounter{logic}%
  \begin{tcolorbox}[
    enhanced,breakable,
    title={Logique \thelogic{} -- #1 \hfill\taglabel{PSEUDO-CODE}},
    colback=bluesoft,
    colframe=blue!55!black,
    coltitle=ink,
    colbacktitle=bluesoft,
    fonttitle=\sffamily\bfseries,
    boxrule=0.35pt,
    arc=0pt,
    left=10pt,right=10pt,top=7pt,bottom=7pt,
    before skip=8pt,after skip=8pt,
    borderline west={3pt}{0pt}{blue}
  ]%
}{\end{tcolorbox}}

\newenvironment{mandatoryexercise}[1]{%
  \refstepcounter{exercise}%
  \par\addvspace{1.05em}
  {\color{line}\hrule height 0.35pt}
  \vspace{0.55em}
  {\sffamily\large\bfseries\color{ink}Exercice \theexercise{} -- #1\par}
  {\taglabel{OBLIGATOIRE}\par}
  \vspace{0.25em}
}{\par\addvspace{0.65em}}

\newenvironment{extensionexercise}[1]{%
  \refstepcounter{extension}%
  \begin{tcolorbox}[
    enhanced,breakable,
    title={Pour aller plus loin \theextension{} -- #1},
    colback=ambersoft,
    colframe=amber!65!black,
    coltitle=ink,
    colbacktitle=ambersoft,
    fonttitle=\sffamily\bfseries,
    boxrule=0.35pt,
    arc=0pt,
    left=10pt,right=10pt,top=8pt,bottom=8pt,
    before skip=8pt,after skip=8pt,
    borderline west={3pt}{0pt}{amber}
  ]%
}{\end{tcolorbox}}

\newenvironment{checkpointbox}[1]{%
  \subsection*{#1}
  \vspace{-0.45em}
}{\vspace{0.15em}}

\newtcolorbox{takeawaybox}[1]{
  enhanced,breakable,
  title={#1},
  colback=greensoft,
  colframe=green,
  coltitle=ink,
  colbacktitle=greensoft,
  fonttitle=\sffamily\bfseries,
  boxrule=0.35pt,
  arc=0pt,
  left=10pt,right=10pt,top=8pt,bottom=8pt,
  before skip=10pt,after skip=8pt,
  borderline west={3pt}{0pt}{green}
}


\begin{document}

\workbooktitle
  {TD2}
  {Boucles, conditions et proportions avec données OWID}
  {Comment utiliser des boucles et des conditions pour classer des observations économiques simples ?}
  {Un fichier \code{td2\_nom\_prenom.py} contenant une boucle, une condition, une proportion, un graphique \code{co2\_td2.png} et une interprétation économique courte.}

\begin{objectivebox}{Objectifs du TD}
\begin{itemize}
  \item Lire une liste de dictionnaires représentant des observations.
  \item Utiliser une boucle \code{for} pour répéter une opération.
  \item Écrire une condition \code{if/elif/else} et compter une proportion.
  \item Construire un graphique simple avec \code{matplotlib} à partir d'une boucle.
  \item Repérer les erreurs fréquentes de boucle, compteur et indentation.
\end{itemize}
\end{objectivebox}

\begin{checkpointbox}{Progression du TD}
Les exercices 1 à 13 sont obligatoires. La notion statistique principale est la proportion. Les seuils utilisés sont pédagogiques : ils servent à apprendre \code{if}, \code{for} et les compteurs, sans aller trop vite vers un programme complet.
\end{checkpointbox}

\begin{checkpointbox}{Source des données}
Les valeurs de CO$_2$ par habitant sont des extraits pédagogiques arrondis issus d'Our World in Data.
\begin{lstlisting}[language=terminal]
https://ourworldindata.org/grapher/co-emissions-per-capita.csv
\end{lstlisting}
\end{checkpointbox}

\begin{pseudocodeexercise}{Classer des observations}
\begin{lstlisting}
# Objectif : classer des pays selon un indicateur
# Entree : liste de pays et valeurs de CO2 par habitant
# Pour chaque pays :
#     lire la valeur
#     comparer la valeur a un seuil
#     afficher une categorie
# Compter combien de pays sont dans chaque categorie
# Sortie : effectifs, proportions, interpretation
\end{lstlisting}
\end{pseudocodeexercise}

\section*{A. Représenter les données}

\begin{mandatoryexercise}{Observer un dictionnaire seul}
Avant de manipuler plusieurs observations, lire une seule observation.
\begin{lstlisting}
obs = {"pays": "France", "co2_pc": 4.7}

print(obs["pays"])
print(obs["co2_pc"])
\end{lstlisting}
Avant d'exécuter, prédire les deux affichages. Modifier ensuite la valeur de \code{co2\_pc} à \code{5.1} et vérifier que seul le résultat affiché change.
\end{mandatoryexercise}

\begin{mandatoryexercise}{Identifier une clé de dictionnaire}
Lire le fragment suivant avant de l'exécuter.
\begin{lstlisting}
obs = {"pays": "France", "co2_pc": 4.7}

print(obs["pays"])
print(obs["population"])
\end{lstlisting}
Prédire quelle ligne fonctionne et quelle ligne produit une erreur. Exécuter ensuite le fragment, noter le type d'erreur en commentaire, puis remplacer \code{"population"} par une clé réellement présente dans le dictionnaire.
\end{mandatoryexercise}

\begin{mandatoryexercise}{Créer une liste de dictionnaires}
Créer la liste suivante. Chaque dictionnaire représente une observation.
\begin{lstlisting}
donnees = [
    {"pays": "France", "co2_pc": 4.7},
    {"pays": "Germany", "co2_pc": 7.9},
    {"pays": "Spain", "co2_pc": 5.0},
    {"pays": "Poland", "co2_pc": 8.4},
    {"pays": "United States", "co2_pc": 14.9},
    {"pays": "India", "co2_pc": 2.0},
    {"pays": "China", "co2_pc": 8.0},
    {"pays": "World", "co2_pc": 4.7}
]
\end{lstlisting}
Afficher le nombre d'observations et la première observation.
\end{mandatoryexercise}

\begin{mandatoryexercise}{Ajouter une observation courte}
Ajouter à la fin de \code{donnees} une observation pour \code{"Brazil"} avec \code{co2\_pc = 2.3}. Relancer ensuite seulement deux contrôles :
\begin{enumerate}
  \item \code{print(len(donnees))} ;
  \item \code{print(donnees[-1])}.
\end{enumerate}
Vérifier que le reste du script pourra utiliser cette nouvelle observation sans modifier les boucles.
\end{mandatoryexercise}

\begin{mandatoryexercise}{Parcourir les observations}
Avec une boucle \code{for}, afficher une phrase par pays :
\begin{lstlisting}
for obs in donnees:
    print(f"{obs['pays']} : {obs['co2_pc']} tonnes par habitant")
\end{lstlisting}
Observer que \code{obs} prend successivement la valeur de chaque dictionnaire de la liste. Modifier ensuite l'affichage pour arrondir à une décimale.
\end{mandatoryexercise}

\section*{B. Conditions et proportions}

\begin{mandatoryexercise}{Classer avec une condition}
On utilise trois catégories pédagogiques :
\begin{itemize}
  \item \code{faible} si \code{co2\_pc < 5} ;
  \item \code{intermediaire} si \code{5 <= co2\_pc < 10} ;
  \item \code{eleve} si \code{co2\_pc >= 10}.
\end{itemize}
Compléter la boucle suivante pour afficher la catégorie de chaque pays.
\begin{lstlisting}
for obs in donnees:
    pays = obs["pays"]
    co2 = obs["co2_pc"]
    if co2 < 5:
        categorie = ...
    elif ...:
        categorie = ...
    else:
        categorie = ...
    print(pays, categorie)
\end{lstlisting}
\end{mandatoryexercise}

\begin{mandatoryexercise}{Tester les seuils avant la boucle complète}
Avant d'appliquer le classement à tous les pays, tester quatre valeurs isolées : \code{4.9}, \code{5.0}, \code{9.9} et \code{10.0}. Pour chacune, écrire en commentaire la catégorie attendue, puis vérifier avec quelques lignes utilisant \code{if/elif/else}.

Expliquer en une phrase pourquoi le cas \code{5.0} doit être classé \code{intermediaire} et non \code{faible}.
\end{mandatoryexercise}

\begin{mandatoryexercise}{Compter une proportion}
Compter le nombre de pays dont les émissions dépassent 5 tonnes par habitant. Calculer ensuite la proportion :
\[
\frac{\text{nombre de pays au-dessus du seuil}}{\text{nombre total de pays}} \times 100.
\]
Compléter d'abord ce squelette :
\begin{lstlisting}
compteur = 0
for obs in donnees:
    if obs["co2_pc"] > 5:
        compteur = compteur + 1

proportion = ...
\end{lstlisting}
Afficher une phrase du type : \code{62.5 \% des observations sont au-dessus du seuil.}
\end{mandatoryexercise}

\begin{pseudocodeexercise}{Préparer des listes pour un graphique}
\begin{lstlisting}
# Objectif : tracer une valeur par pays
# Entree : liste de dictionnaires
# 1. Creer une liste vide pour les noms de pays
# 2. Creer une liste vide pour les valeurs de CO2
# 3. Parcourir les observations une par une
# 4. Ajouter le pays et la valeur dans les deux listes
# 5. Tracer puis sauvegarder un graphique
# Sortie : figure png
\end{lstlisting}
\end{pseudocodeexercise}

\begin{mandatoryexercise}{Graphique des émissions par habitant}
Construire deux listes avec une boucle, puis tracer un graphique en barres.
\begin{lstlisting}
import matplotlib.pyplot as plt

noms_pays = []
valeurs_co2 = []

for obs in donnees:
    noms_pays.append(obs["pays"])
    valeurs_co2.append(obs["co2_pc"])

plt.figure()
plt.bar(noms_pays, valeurs_co2)
plt.ylabel("Tonnes de CO2 par habitant")
plt.title("Emissions de CO2 par habitant - extrait OWID")
plt.xticks(rotation=45, ha="right")
plt.tight_layout()
plt.savefig("co2_td2.png", dpi=300)
\end{lstlisting}
\end{mandatoryexercise}

\begin{mandatoryexercise}{Créer une fonction de classement}
Compléter la fonction \code{classe\_co2(valeur)}. Elle doit renvoyer \code{"faible"}, \code{"intermediaire"} ou \code{"eleve"}.
\begin{lstlisting}
def classe_co2(valeur):
    if valeur < 5:
        return ...
    elif valeur < 10:
        return ...
    else:
        return ...
\end{lstlisting}
La tester d'abord avec \code{classe\_co2(4.7)}, \code{classe\_co2(7.9)} et \code{classe\_co2(14.9)}, puis l'utiliser dans une boucle pour afficher le classement de tous les pays.
\end{mandatoryexercise}

\begin{mandatoryexercise}{Débogage : boucle, compteur, indentation}
Les trois fragments suivants contiennent une erreur fréquente. Pour chacun, identifier le type d'erreur probable ou le résultat logique faux, puis proposer une correction.
\begin{lstlisting}
# Fragment A
for obs in donnees:
print(obs["pays"])

# Fragment B
compteur = compteur + 1

# Fragment C
compteur = 0
for obs in donnees:
    if obs["co2_pc"] > 5:
    compteur = compteur + 1
\end{lstlisting}
Ne pas intégrer ces fragments fautifs au script final.
\end{mandatoryexercise}

\section*{C. Interpréter}

\begin{mandatoryexercise}{Question économique}
Ajouter quatre commentaires à la fin du script :
\begin{enumerate}
  \item quel pays de l'extrait a la valeur la plus élevée ?
  \item la France est-elle au-dessus ou au-dessous de la valeur mondiale ?
  \item quelle proportion de l'extrait dépasse 5 tonnes par habitant ?
  \item pourquoi ne faut-il pas confondre CO$_2$ par habitant et CO$_2$ total ?
\end{enumerate}
\end{mandatoryexercise}

\begin{mandatoryexercise}{Contrôle final}
Le script doit contenir au minimum : une liste de dictionnaires, une boucle \code{for}, une condition \code{if/elif/else}, un compteur, une proportion, une fonction, un graphique \code{matplotlib} et une interprétation.
\end{mandatoryexercise}

\section*{D. Entraînement facultatif}

\begin{extensionexercise}{Dix petits entraînements}
Choisir au moins deux questions si le parcours obligatoire est terminé.
\begin{enumerate}
  \item Changer le seuil de 5 à 7 tonnes et recalculer la proportion.
  \item Créer une liste contenant seulement les pays classés \code{eleve}.
  \item Calculer la moyenne de \code{co2\_pc} avec une boucle, sans \code{sum()}.
  \item Ajouter un pays et vérifier que le code fonctionne encore.
  \item Afficher les pays au-dessous de la valeur mondiale.
  \item Compter séparément les pays \code{faible}, \code{intermediaire} et \code{eleve}.
  \item Écrire une fonction \code{au\_dessus\_seuil(valeur, seuil)}.
  \item Trier manuellement deux pays : lequel a la valeur la plus élevée ?
  \item Écrire une phrase sur la limite d'un classement par seuil.
  \item Expliquer pourquoi une proportion dépend de l'échantillon choisi.
\end{enumerate}
\end{extensionexercise}

\begin{takeawaybox}{À retenir}
\begin{itemize}
  \item Une boucle \code{for} permet d'appliquer la même logique à chaque observation.
  \item Une condition transforme une comparaison booléenne en décision.
  \item Un compteur doit être initialisé avant la boucle, puis modifié au bon niveau d'indentation.
  \item Un graphique peut se préparer en construisant explicitement les listes à afficher.
  \item Une proportion dépend toujours de l'ensemble d'observations choisi.
\end{itemize}
\end{takeawaybox}

\end{document}
